% tests/zebra-test.tex
%
% Shared assertion helpers for zebra's behavioural regression tests.
%
% Rather than diffing golden log files (fragile across TeX
% distributions and kernel versions), each test asserts the *semantic*
% outcome of the dedup machinery: how many distinct notes were actually
% allocated.  \g__zebra_note_id_int counts one per real allocation
% (Case 4 in \@@_note_numbered:nnn); sbox reuse, stable-key suppression
% and content-signature suppression do NOT bump it.  The per-type
% counter g__zebra_note_count_<type>_int behaves the same way.  So these
% two numbers are exactly what a dedup regression (over- or
% under-collapsing) would move.
%
% A failed assertion raises a LaTeX error, which the test runner
% (tests/runtests.lua, and the legacy Makefile target) already treats as
% a failure -- this reuses the existing \PackageError pattern that a few
% tests spelt out by hand.  The error message reports got/want.
%
% Counts only converge once the .aux signatures from the previous pass
% are read back, so call these at end of document and let the runner
% read the last of its three pdflatex passes.
\ExplSyntaxOn
\msg_new:nnn { zebra-test } { count-mismatch }
  { Behavioural~regression:~#1. }
\cs_new_protected:Npn \__zebra_test_assert:nnn #1#2#3
  {
    \int_compare:nNnF {#2} = {#3}
      {
        \msg_error:nnx { zebra-test } { count-mismatch }
          { #1~got~\int_eval:n {#2},~want~#3 }
      }
  }
% \ZebraAssertTotal{<n>}: total distinct allocations across all types.
\NewDocumentCommand \ZebraAssertTotal { m }
  { \__zebra_test_assert:nnn { total } { \g__zebra_note_id_int } {#1} }
% \ZebraAssertCount{<type>}{<n>}: distinct notes of <type>.
\NewDocumentCommand \ZebraAssertCount { m m }
  {
    \__zebra_test_assert:nnn { count[#1] }
      { \int_use:c { g__zebra_note_count_#1_int } } {#2}
  }
\ExplSyntaxOff
